%==============================================================================
% Sjabloon onderzoeksvoorstel bachproef
%==============================================================================
% Gebaseerd op document class `hogent-article'
% zie <https://github.com/HoGentTIN/latex-hogent-article>

% Voor een voorstel in het Engels: voeg de documentclass-optie [english] toe.
% Let op: kan enkel na toestemming van de bachelorproefcoördinator!
\documentclass{hogent-article}

% Invoegen bibliografiebestand
\addbibresource{voorstel.bib}

% Informatie over de opleiding, het vak en soort opdracht
\studyprogramme{Professionele bachelor toegepaste informatica}
\course{Bachelorproef}
\assignmenttype{Onderzoeksvoorstel}
% Voor een voorstel in het Engels, haal de volgende 3 regels uit commentaar
% \studyprogramme{Bachelor of applied information technology}
% \course{Bachelor thesis}
% \assignmenttype{Research proposal}

\academicyear{2025-2026}

\title{De automatisatiestack van de toekomst op z/OS met Rexx en Python}

\author{Vincent Gielen}
\email{vincent.gielen@student.hogent.be}

% TODO: Geef de co-promotor op
\supervisor[Co-promotor]{B. Tjassens (Rabobank, \href{mailto:	arthurcoucke@gmail.com}{arthurcoucke@gmail.com})}

% Binnen welke specialisatierichting uit 3TI situeert dit onderzoek zich?
% Kies uit deze lijst:
%
% - Mobile \& Enterprise development
% - AI \& Data Engineering
% - Functional \& Business Analysis
% - System \& Network Administrator
% - Mainframe Expert
% - Als het onderzoek niet past binnen een van deze domeinen specifieer je deze
%   zelf
%
\specialisation{Mainframe Expert}
\keywords{Mainframe, Modernisatie, Rexx, Python, JCL, vergelijkende studie}

\begin{document}

\begin{abstract}
  Dit onderzoek wil de toekomst van de automatisatiestack op z/OS (het voornaamste besturingssysteem op mainframe) onderzoeken. Het focust daarbij voornamelijk op Rexx en Python, twee mogelijke kandidaten voor de automatisatiestack op z/OS met elk hun voor- en nadelen. Rexx draagt bij tot de stabiliteit en compatibiliteit van mainframes, terwijl Python dan weer meer gepast lijkt voor de snelle vernieuwing van deze tijd. Wanneer er voor een bepaalde taak moet gekozen worden voor een van die twee scripting talen, is het belangrijk om een goed overzicht te hebben van alle voor- en nadelen. \textcolor{red} {Gezien de modernisatiegolf die momenteel aan de gang is, en het gebrek aan kennis van Rexx bij velen op de arbeidsmarkt, lijkt Python de uitgelezen kandidaat. Maar is dat wel zo?} Dit onderzoek tracht volgende onderzoeksvraag te beantwoorden: Hoe ziet de automatisatie stack van de toekomst eruit op z/OS en wat staat er in de weg om van Python de nieuwe standaard te maken?
  
  \qquad Om de onderzoeksvraag te beantwoorden wordt eerst een literatuurstudie uitgevoerd. Op basis daarvan volgt een algemene vergelijkende studie van Rexx en Python, waar wordt getracht om enkele veel gebruikte Rexx commando’s te vervangen door Python. Tot slot wordt er met behulp van een PoC dieper ingegaan op één specifieke use case van Rexx, namelijk de File Tailoring Services voor het genereren van Job Control Language (JCL).
  
  \qquad Vermoedelijk wijzen de literatuurstudie en de vergelijkende studie uit dat, enerzijds, Python nodig is om de toekomst van mainframe te verzekeren, omdat het de instapdrempel verlaagt voor jongere mensen en veel flexibiliteit geeft. Anderzijds is de efficiëntie van Rexx en de specifieke bouw van Rexx voor scripting en mainframes vermoedelijk niet te evenaren door Python. Uit de PoC zal waarschijnlijk opnieuw blijken dat hoewel Python enorme flexibiliteit biedt, de specifieke bouw van Rexx ervoor zorgt dat alles efficiënter verloopt .
  
\end{abstract}

\tableofcontents

% De hoofdtekst van het voorstel zit in een apart bestand, zodat het makkelijk
% kan opgenomen worden in de bijlagen van de bachelorproef zelf.
%---------- Inleiding ---------------------------------------------------------

\section{Inleiding}%
\label{sec:inleiding}

De mainframe is een van de oudste technologieën, en nog steeds een van de meeste gebruikte. Het merendeel van de grootste bedrijven in zowat elke sector gebruiken een mainframe (bank, retail, logistiek,...). Dat de technologie oud is betekent echter niet dat die niet goed werkt. Als het gaat om hoeveelheid transacties is de mainframe veruit de meest efficiënte, betrouwbare en veilige technologie. Door een combinatie van redenen, hebben mainframes al jarenlang een reputatie van een ter dood veroordeelde technologie. Dit heeft ervoor gezorgd dat er in een hele generatie IT'ers zich slechts weinigen hebben gespecialiseerd in Mainframe. Als gevolg krampen bedrijven nu met een groot tekort aan beschikbare werkkrachten die overweg kunnen met de unieke technologie achter mainframes.

Een van die unieke technologieën is Rexx. Rexx is een scripting taal, vooral gebruikt voor automatisatie. Een vergelijkbare, modernere technologie is Python. Vandaar dat Python vaak wordt gebruikt ter vervanging van Rexx. 

Een van de vele taken waarvoor zowel Rexx als Python voor gebruikt worden is het creëren van JCL (Job Control Language) bestanden met de hulp van templates. Bij Rexx kan gebruikt gemaakt worden van ISPF File Tailoring, terwijl bij Python Jinja kan gebruikt worden. 

%---------- Stand van zaken ---------------------------------------------------

\section{Literatuurstudie}%
\label{sec:literatuurstudie}

Hier beschrijf je de \emph{state-of-the-art} rondom je gekozen onderzoeksdomein, d.w.z.\ een inleidende, doorlopende tekst over het onderzoeksdomein van je bachelorproef. Je steunt daarbij heel sterk op de professionele \emph{vakliteratuur}, en niet zozeer op populariserende teksten voor een breed publiek. Wat is de huidige stand van zaken in dit domein, en wat zijn nog eventuele open vragen (die misschien de aanleiding waren tot je onderzoeksvraag!)?

Je mag de titel van deze sectie ook aanpassen (literatuurstudie, stand van zaken, enz.). Zijn er al gelijkaardige onderzoeken gevoerd? Wat concluderen ze? Wat is het verschil met jouw onderzoek?

Verwijs bij elke introductie van een term of bewering over het domein naar de vakliteratuur, bijvoorbeeld~\autocite{Prechelt2000}! Denk zeker goed na welke werken je refereert en waarom.

Draag zorg voor correcte literatuurverwijzingen! \autocite{SagersEtAl2018} Een bronvermelding \textcite{SochovaKunce2012} hoort thuis \emph{binnen} de zin waar je je op die bron baseert, dus niet er buiten! Maak meteen een verwijzing als je \autocite{MurphyEtAl2010} gebruik maakt van een bron. Doe dit dus \emph{niet} aan het einde van een lange paragraaf. Baseer nooit teveel aansluitende tekst op eenzelfde bron.

Als je informatie over bronnen verzamelt in JabRef, zorg er dan voor dat alle nodige info aanwezig is om de bron terug te vinden (zoals uitvoerig besproken in de lessen Research Methods).

% Voor literatuurverwijzingen zijn er twee belangrijke commando's:
% \autocite{KEY} => (Auteur, jaartal) Gebruik dit als de naam van de auteur
%   geen onderdeel is van de zin.
% \textcite{KEY} => Auteur (jaartal)  Gebruik dit als de auteursnaam wel een
%   functie heeft in de zin (bv. ``Uit onderzoek door Doll & Hill (1954) bleek
%   ...'')

Je mag deze sectie nog verder onderverdelen in subsecties als dit de structuur van de tekst kan verduidelijken.

%---------- Methodologie ------------------------------------------------------
\section{Methodologie}%
\label{sec:methodologie}

Hier beschrijf je hoe je van plan bent het onderzoek te voeren. Welke onderzoekstechniek ga je toepassen om elk van je onderzoeksvragen te beantwoorden? Gebruik je hiervoor literatuurstudie, interviews met belanghebbenden (bv.~voor requirements-analyse), experimenten, simulaties, vergelijkende studie, risico-analyse, PoC, \ldots?

Valt je onderwerp onder één van de typische soorten bachelorproeven die besproken zijn in de lessen Research Methods (bv.\ vergelijkende studie of risico-analyse)? Zorg er dan ook voor dat we duidelijk de verschillende stappen terug vinden die we verwachten in dit soort onderzoek!

Vermijd onderzoekstechnieken die geen objectieve, meetbare resultaten kunnen opleveren. Enquêtes, bijvoorbeeld, zijn voor een bachelorproef informatica meestal \textbf{niet geschikt}. De antwoorden zijn eerder meningen dan feiten en in de praktijk blijkt het ook bijzonder moeilijk om voldoende respondenten te vinden. Studenten die een enquête willen voeren, hebben meestal ook geen goede definitie van de populatie, waardoor ook niet kan aangetoond worden dat eventuele resultaten representatief zijn.

Uit dit onderdeel moet duidelijk naar voor komen dat je bachelorproef ook technisch voldoen\-de diepgang zal bevatten. Het zou niet kloppen als een bachelorproef informatica ook door bv.\ een student marketing zou kunnen uitgevoerd worden.

Je beschrijft ook al welke tools (hardware, software, diensten, \ldots) je denkt hiervoor te gebruiken of te ontwikkelen.

Probeer ook een tijdschatting te maken. Hoe lang zal je met elke fase van je onderzoek bezig zijn en wat zijn de concrete \emph{deliverables} in elke fase?

%---------- Verwachte resultaten ----------------------------------------------
\section{Verwacht resultaat, conclusie}%
\label{sec:verwachte_resultaten}

Hier beschrijf je welke resultaten je verwacht. Als je metingen en simulaties uitvoert, kan je hier al mock-ups maken van de grafieken samen met de verwachte conclusies. Benoem zeker al je assen en de onderdelen van de grafiek die je gaat gebruiken. Dit zorgt ervoor dat je concreet weet welk soort data je moet verzamelen en hoe je die moet meten.

Wat heeft de doelgroep van je onderzoek aan het resultaat? Op welke manier zorgt jouw bachelorproef voor een meerwaarde?

Hier beschrijf je wat je verwacht uit je onderzoek, met de motivatie waarom. Het is \textbf{niet} erg indien uit je onderzoek andere resultaten en conclusies vloeien dan dat je hier beschrijft: het is dan juist interessant om te onderzoeken waarom jouw hypothesen niet overeenkomen met de resultaten.



\printbibliography[heading=bibintoc]

\end{document}